% rosetta/appendix-sym.tex — 附录 C：三张速查表（左右约定、符号索引、状态分级）
\section{速查表}
\label{app:sym}

本附录只做三件事：表~C.1 固定左右与模型约定，表~C.2 给出符号的首现位置，表~C.3 标明重要结论的证据等级。读正文遇到「这个符号在哪一节」「这句话是定理还是猜测」「$J$ 侧的膜到底对偶到哪一侧」时，先查这里。三张表都用 \texttt{longtable}，可以跨页；表里只放短句，较长的说明放在每张表后的注里。

\subsection{表 C.1：左右约定与四组容易弄反的配对}

先读这张表，再读正文里任何一句带「对偶」「镜像」「特征」的话。四组配对指电与磁、零膜与平坦丛、$J$ 侧与 $K$ 侧、$B_{cc}$ 与 $A_{cc}$；第一行与第二行先固定「谁在哪一侧」这个大前提。

\begingroup
\footnotesize
\renewcommand{\arraystretch}{1.16}
\begin{longtable}{@{}p{3.0cm}p{2.9cm}p{4.8cm}p{3.85cm}@{}}
\toprule
\bfseries 对象 & \bfseries 侧／模型 & \bfseries 正确读法 & \bfseries 容易反读成什么 \\
\midrule
\endfirsthead
\toprule
\bfseries 对象 & \bfseries 侧／模型 & \bfseries 正确读法 & \bfseries 容易反读成什么 \\
\midrule
\endhead
\bottomrule
\endlastfoot

$G$ 与 ${}^LG$ 住在哪一侧 & $G$：$\mathrm{Bun}_G$ 与 $\mathcal M(G,C)$；${}^LG$：$\mathrm{Loc}_{{}^LG}(C)$ 与 $\mathcal M_H({}^LG,C)$ & 几何 Langlands 左侧是 $G$-丛上的 $\mathcal D$-模，右侧是 ${}^LG$-局部系统上的相干层；$S$-对偶交换 $G$ 与 ${}^LG$ & 以为两侧都是 $G$；或以为 $\mathrm{Bun}_G$ 的对偶是 $\mathrm{Bun}_{{}^LG}$，其实是 $\mathrm{Loc}_{{}^LG}$ \\
\addlinespace[3pt]
膜与它的「数学翻译」 & 膜：$\mathcal M_H$ 上的边界条件；翻译：$\mathrm{Bun}_G$ 上的层或模 & 电特征膜（零膜）与 Hecke 特征层是同一对象在两个范畴里的化身，中间隔着镜像对称与 $B_{cc}$ 翻译 & 直接把膜名当层名，说「电特征膜就是 $\mathrm{Bun}_G$ 上的 Hecke 特征层」 \\
\addlinespace[3pt]
电特征膜与磁特征膜 & 电：$\Psi=\infty$ 的零膜；磁：$\Psi=0$ 的 type $\mathcal F$ 膜，支撑在 Hitchin 纤维上 & 两者是同一 Hitchin 模空间上的两类膜，$S$ 交换之；本征值分别是 $\mathcal E_p(R)$ 的纤维与 ${}^LR$ 的平坦丛 & 以为电膜在 $\mathrm{Bun}_G$ 侧、磁膜在 $\mathrm{Loc}$ 侧；或把二者当成同一块膜的两个名字 \\
\addlinespace[3pt]
零膜 $\mathcal O_E$ 与平坦 ${}^LG$-丛 $E$ & $E\in\mathrm{Loc}_{{}^LG}(C)$，对应 $\mathcal M_H({}^LG,C)$ 的 $J$ 侧；$\mathcal O_E$ 是 $E$ 处的摩天层 & $\mathcal O_E$ 的镜像是对偶 Hitchin 纤维上的 $A$-膜 $\mathcal F_E$；再经 $B_{cc}$ 翻译才落到 $\mathrm{Bun}_G$ 上 & 把 $E$ 读成 $G$-丛；把 $\mathcal O_E$ 读成 $\mathrm{Bun}_G$ 上的结构层 \\
\addlinespace[3pt]
$J$ 侧 $B$-膜与 $K$ 侧 $A$-膜 & $B$-模型在 $\mathcal M_H({}^LG,C)$ 的复结构 $J$；$A$-模型在 $\mathcal M_H(G,C)$ 的辛形式 $\omega_K$ & 同调镜像对称把它们配成一对；$K$ 与 $J$ 是两个不同的复结构，$A$-模型只依赖 $\omega_K$ & 以为 $B$-膜配 $B$-膜；或把 $\omega_K$ 说成 $J$ 的 Kähler 形式 \\
\addlinespace[3pt]
$B_{cc}$ 与文献中的 $A_{cc}$ & $\mathcal M_H(G,C)$ 上的典范余各向同性 $A$-膜 & 同一对象：KW06 写 $\mathcal B_{c.c.}$，综述与词典正文写 $A_{cc}$；$cc$ 是 canonical coisotropic 的缩写 & 以为是两块膜，一块 $A$、一块 $B$；或以为 $B_{cc}$ 是 $B$-膜 \\

\end{longtable}
\endgroup

注：第一行问「谁在哪一侧」，第三、四、五行问「膜在哪一侧」，第二行问「膜与它的翻译是否同一范畴」。最容易错的是第二行：电特征膜是 $\mathcal M_H$ 上的零膜，Hecke 特征层是 $\mathrm{Bun}_G$ 上的模，二者只通过镜像对称与 \S11 的 $B_{cc}$ 函子相连。第六行是记号冲突：项目正文里 $A_{cc}$ 与 $B_{cc}$ 同时出现，本表统一读作 $B_{cc}$；它绝不是「$B$-膜」，而是典范余各向同性（canonical coisotropic）$A$-膜。

\subsection{表 C.2：符号—首现位置索引}

第二列给含义，第三列给首现的节与印刷页（印刷页 $=$ 页脚数字；文献页码按「节 页」写，例如 \S4.1 p.45）。符号按主题分组，不是按字母序。

\begingroup
\footnotesize
\renewcommand{\arraystretch}{1.14}
\begin{longtable}{@{}p{3.4cm}p{7.4cm}p{4.3cm}@{}}
\toprule
\bfseries 符号 & \bfseries 含义 & \bfseries 首现节与印刷页 \\
\midrule
\endfirsthead
\toprule
\bfseries 符号 & \bfseries 含义 & \bfseries 首现节与印刷页 \\
\midrule
\endhead
\bottomrule
\endlastfoot

\multicolumn{3}{@{}l}{\bfseries 一、群、格与对偶}\\
\addlinespace[2pt]
$G$ & 规范群；四维理论的规范群，也决定二维靶空间 & \S1 p.2 \\
${}^LG$ & GNO 对偶群，即 Langlands 对偶群 & \S1 p.3 \\
$\hat G$ & 对偶群的另一种写法；原文一律用 ${}^LG$ & \S?（原文未用此记号） \\
$\mathfrak g$ & $G$ 的李代数 & \S2.1 p.6 \\
$n_{\mathfrak g}$ & lacing 数：单连通为 1，$B_n,C_n,F_4$ 为 2，$G_2$ 为 3 & \S2.2 p.11 \\
$\Lambda$ & 根格（原文如此） & \S7.1 p.103 \\
$\Lambda^\vee$ & 权格（原文如此，与常见记号相反） & \S7.1 p.103 \\
$Z(G)$ & 中心；它分类 $S^2$ 上的 $G$-丛 & \S6.2 p.84 \\
\addlinespace[4pt]
\multicolumn{3}{@{}l}{\bfseries 二、耦合、扭化与场}\\
\addlinespace[2pt]
$\tau$ & 复耦合 $\theta/2\pi+4\pi i/e^2$ & \S2.2 p.11 \\
$\theta$ & $\theta$ 角 & \S1 p.3 \\
$e$ & 规范耦合常数 & \S2.2 p.11 \\
$\Psi$ & 典范参数；$\Psi=\infty$ 取电、$\Psi=0$ 取磁 & \S2.2 p.11（定义 \S3.5 p.35） \\
$t$ & 扭化参数 $t=v/u$；$t=\pm1$ 给 $A$-模型，$t=\pm i$ 给 $B$-模型 & \S3.1 p.21 \\
$Q$ & 扭化标量超荷 $Q=uQ_\ell+vQ_r$ & \S3.1 p.18 \\
$Q_\ell,Q_r$ & 两个反交换超荷，扭化取它们的组合 & \S3.1 p.18 \\
$A$ & 四维规范场（联络 $A_\mu$） & \S2.1 p.6 \\
$\phi$ & 四维 $\mathcal N=4$ 的标量场 & \S2.1 p.6 \\
$\varphi$ & 二维 Higgs 场，$K_C\otimes\mathrm{ad}E$ 的截面 & \S4.2 p.49 \\
$\sigma$ & 标量组合 $(\phi_4-i\phi_5)/\sqrt2$ & \S3.1 p.17 \\
$\epsilon$ & 超对称参数（$\mathcal S^+$ 中的常数旋量） & \S2.1 p.9 \\
$\delta$ & de Rham 微分，超荷在形式上的作用 & \S4.1 p.45 \\
$\hat\phi$ & Coulomb 分支的标量场（中心荷相位） & \S2.2 p.14 \\
\addlinespace[4pt]
\multicolumn{3}{@{}l}{\bfseries 三、Hitchin 模空间与可积系统}\\
\addlinespace[2pt]
$\mathcal M_H(G,C)$ & $G$ 的 Hitchin 模空间（超凯勒） & \S1 p.2（定义 \S4.1 p.45） \\
$\mathcal M_H({}^LG,C)$ & 对偶 Hitchin 模空间，镜像的靶 & \S5.4 p.71 \\
$I,J,K$ & 三个超凯勒复结构 & \S4.1 p.42 \\
$\omega_I,\omega_J,\omega_K$ & 三个 Kähler 形式 & \S4.1 p.44 \\
$\Omega_I,\Omega_J,\Omega_K$ & 三个全纯辛形式 & \S4.1 p.44 \\
$B$ & Hitchin 基（Casimir 特征多项式空间） & \S4.3 p.53 \\
$\pi$ & Hitchin 纤维化 $\pi:\mathcal M_H\to B$ & \S4.3 p.53 \\
$\mathcal F$ & Hitchin 纤维；type $\mathcal F$ 膜的支撑 & \S5.5 p.74 \\
$H_R$ & Hitchin 哈密顿量（原文写 $H_\alpha$），纤维由 $H_\alpha=h_\alpha$ 定义 & \S5.5 p.74 \\
$h_\alpha$ & 哈密顿量在纤维上的取值 & \S5.5 p.74 \\
\addlinespace[4pt]
\multicolumn{3}{@{}l}{\bfseries 四、模空间、层与翻译}\\
\addlinespace[2pt]
$\mathcal M(G,C)$ & $G$-丛模空间 & \S7.1 p.103 \\
$\mathrm{Bun}_G$ & $G$-丛的栈（几何 Langlands 左侧） & \S10.1 p.148 \\
$\mathrm{Loc}_{{}^LG}(C)$ & 平坦 ${}^LG$-丛模空间；综述与词典记号 & \S?（原文用 $\mathcal Y$，\S3.2 p.25） \\
$\mathcal Y({}^LG,C)$ & 同态 $\pi_1(C)\to{}^LG$ 的模空间 & \S3.2 p.25 \\
$\mathcal D$ & 微分算子层，$\mathcal D$-模的来源 & \S1 p.2（\S11.2 p.196） \\
$\mathcal D_{K^{1/2}}$ & 扭化微分算子层（差一个 $K^{1/2}$） & \S11.2 p.196 \\
$K_C$ & 典范线丛 & \S4.2 p.49 \\
$\mathcal O_E$ & $E$ 处的摩天层（零膜的 Chan-Paton 层） & \S8.2 p.121 \\
$F_E$ & $E$ 上的 Hecke 特征层；词典记号 & \S?（原文称 Hecke eigensheaf，\S6.4 p.95） \\
$\mathcal F_E$ & 对偶 Hitchin 纤维上的 $A$-膜 & \S5.5 p.74 \\
$V_R$ & 表示 $R$ 的向量空间（特征膜的本征值） & \S6.4 p.95 \\
$\mathcal B$ & 膜（二维模型的边界条件） & \S5.2 p.65 \\
$\mathcal B_p$ & 支撑在点 $p$ 的零膜 & \S5.5 p.74 \\
$U_x$ & 支撑在 $x$ 的摩天层，零膜的 Chan-Paton 层 & \S8.2 p.121 \\
$\mathcal L$ & 线丛（Chan-Paton 平坦线丛） & \S5.5 p.74 \\
\addlinespace[4pt]
\multicolumn{3}{@{}l}{\bfseries 五、算子与特殊记号}\\
\addlinespace[2pt]
$W_q(R)$ & Wilson 线算子：表示 $R$、点 $q$ & \S6.1 p.79 \\
$T({}^LR,S)$ & 't Hooft 线算子：${}^LR$ 标号，沿 $S$ & \S6.1 p.79（标号 \S6.2 p.84） \\
$\rho$ & $U(1)\to G$ 同态，给出 't Hooft 奇异类 & \S6.2 p.84 \\
$E_p(R)$ & 泛丛在 $p$ 的纤维（表示 $R$），Wilson 算子的本征值 & \S8.2 p.121 \\
$T({}^Lw,z_0)$ & 带参数 $({}^Lw,z_0)$ 的 't Hooft 算子（深水区） & \S10.5 p.179 \\
$\bar{\mathcal Y}({}^Lw;z_0)$ & 奇异 Bogomolny 解的修正空间 & \S10.5 p.182 \\
$\mathcal M^{\mathrm{mon}}_k$ & $k$ 个单极子的模空间；词典记号 & \S10.2 p.155 \\
$\bar M_k$ & 单极子模空间的紧化；词典记号 & \S10.2 p.155 \\
$B_{cc}$ & 典范余各向同性 $A$-膜（原文 $\mathcal B_{c.c.}$） & \S11.1 p.185 \\
$A_{cc}$ & 同一对象，综述与词典正文的记号 & \S?（同一对象，见 \S11.1 p.185） \\
$c$ & $cc$ 即 canonical coisotropic 的缩写 & \S11.1 p.185 \\

\end{longtable}
\endgroup

注：印刷页与文献页码的换算是「页脚数字 $=$ 印刷页」，表中已按这个约定写出；原文没有的数字（例如 \S3.1 内 $t$ 的确切页）由软拷贝匹配得到，误差约一页。标 \S? 的条目共四条：$\hat G$、$\mathrm{Loc}_{{}^LG}$、$F_E$、$A_{cc}$，另有 $\mathcal M^{\mathrm{mon}}_k$ 与 $\bar M_k$ 是词典自造的记号。原文的写法与本表略有出入：$B_{cc}$ 原文记 $\mathcal B_{c.c.}$，$E_p(R)$ 原文记 $\mathcal E_q(R)$，$\mathrm{Loc}$ 原文用 $\mathcal Y$（de Rham 版）并只在 \S10.4 p.171 说 locally constant。

\subsection{表 C.3：本论文各部分的状态：定理 / 推导 / 纲要}

第三列只有三种取值：\textbf{定理（引用）} 指原文引用他人已证的结论；\textbf{原文推导（argued）} 指原文给出论证但未写成严格证明；\textbf{纲要（原文只说方向）} 指原文只说「应该如此」或把它留给未来。查一个结论能不能当定理用，先看这一列。

\begingroup
\footnotesize
\renewcommand{\arraystretch}{1.14}
\begin{longtable}{@{}p{4.5cm}p{2.9cm}p{3.2cm}p{4.3cm}@{}}
\toprule
\bfseries 结论 & \bfseries 位置 & \bfseries 状态 & \bfseries 说明 \\
\midrule
\endfirsthead
\toprule
\bfseries 结论 & \bfseries 位置 & \bfseries 状态 & \bfseries 说明 \\
\midrule
\endhead
\bottomrule
\endlastfoot

$S$-对偶给出二维镜像对称 & \S5.3--\S5.4 p.70--74 & 原文推导（argued） & 靠 Hitchin 纤维环面的 SYZ 对偶论证，未给严格证明 \\
\addlinespace[3pt]
镜像不依赖 $C$ 的复结构 & \S5.5 p.74 & 原文推导（argued） & 由四维扭化 TQFT 的拓扑性推出 \\
\addlinespace[3pt]
电、磁特征膜在 $S$ 下互换 & \S6.4 p.95--97 & 原文推导（argued） & $S$ 把 $\Psi$ 送到 $-1/n_{\mathfrak g}\Psi$，交换两类线算子 \\
\addlinespace[3pt]
't Hooft 算子的 OPE 是向量空间 & \S10.4 p.167 & 原文推导（argued） & 与仿射 Grassmann 的卷积代数逐项对照 \\
\addlinespace[3pt]
单极子模空间等于 Hecke 修正空间 & \S9.1--\S9.2 p.128--135 & 原文推导（argued） & 逐点比较冒泡构型与修正的定义 \\
\addlinespace[3pt]
单极子空间同构于 Schubert 簇 & \S9.3 p.147 & 定理（引用） & 引用仿射 Grassmann 的已知结论 \\
\addlinespace[3pt]
$B_{cc}$ 把 $K$ 上的 $A$-膜送到 $\mathcal D$-模 & \S11.1--\S11.2 p.185--196 & 原文推导（argued） & 只用「c.c. 弦无瞬子修正」的论证 \\
\addlinespace[3pt]
$B_{cc}$ 在单连通 $G$ 时唯一 & \S11.1 p.185 & 原文推导（argued） & 差一个平坦线丛；非单连通时变多 \\
\addlinespace[3pt]
$A$-膜范畴等价于扭化余切丛的 $B$-膜范畴 & \S11.3 p.201 & 定理（引用） & 引前人关于全纯辛流形量子化的结论 \\
\addlinespace[3pt]
$\mathcal D$-模对应相干层（差 $K^{1/2}$） & \S11.2--\S11.3 p.196--201 & 定理（引用） & 引前人等价；$K^{1/2}$ 的细节原文从略 \\
\addlinespace[3pt]
几何 Langlands 等价于 $S$-对偶 & \S1 p.2、\S5.4 p.74 & 纲要（原文只说方向） & 作为猜想级陈述给出，未证明 \\
\addlinespace[3pt]
零膜对应电边界条件（四维提升） & \S12.2 p.211 & 原文推导（argued） & 同一份 $(B,B,B)$ 数据提升不唯一 \\
\addlinespace[3pt]
$(B,A,A)$ 膜对应磁边界条件 & \S12.3 p.215 & 原文推导（argued） & 由镜像的 Hitchin 纤维给出 \\
\addlinespace[3pt]
$(A,B,A)$ 膜对应 c.c. 边界条件 & \S12.4 p.217 & 原文推导（argued） & c.c. 膜的边界可观察量由时间反演论证 \\
\addlinespace[3pt]
边界可观察量无量子修正 & \S12.4 p.217 & 原文推导（argued） & 只用时间反演与超对称，不是逐阶计算 \\
\addlinespace[3pt]
只处理非分歧情形，分歧情形只给方向 & \S1 p.5（不在 \S12） & 纲要（原文只说方向） & 原文仅在引言声明，\S12 未讨论分歧 \\

\end{longtable}
\endgroup

注：\S11 与 \S12 的绝大多数条目是「原文推导（argued）」，这正是两轮盲测报告说「能复述主线、不能重建推导」的原文根据；把它们当定理引用会高估原文。另需更正一处常见误记：非分歧与分歧的边界不在 \S12。全文只有 \S1 p.5 一句声明「只考虑非分歧情形，分歧情形想必有类似想法」，\S12 讨论的是把二维膜提升回四维边界条件，与分歧无关。
